% from ega2-1-fr.tex
\oldpage[II]{190}
est dense dans $V$, donc $h_0$ est un \emph{homéomorphisme} de $U$ sur $V$.
Pour démontrer le corollaire, il suffira de voir que
$\lambda:\mathcal{O}_V\to(h_0)_*(\mathcal{O}_U)$ est un isomorphisme~: on
pourra appliquer alors (\hyperref[II.8.11.1-fr]{8.11.1}) qui prouvera que
l'application (\hyperref[II.8.11.1.1-fr]{8.11.1.1}) est bijective (les fibres
$h_0^{-1}(x)$ étant chacune réduite à un seul point)~; donc $h$ sera un
isomorphisme. La question étant évidemment locale sur $V$, on peut supposer
que $V=\operatorname{Spec}(A)$ est affine d'anneau intègre et intégralement
clos (\hyperref[II.8.8.6.1-fr]{8.8.6.1})~; $h$ correspond alors
(\hyperref[I.2.2.4-fr]{I, 2.2.4}) à un homomorphisme
$\varphi:A\to\Gamma(U,\mathcal{O}_U)$ et tout revient à voir que $\varphi$
est un isomorphisme. Or, si $K$ est le corps des fractions de $A$,
$\Gamma(U,\mathcal{O}_U)$ a par hypothèse $K$ pour corps des fractions et $A$
est un sous-anneau de $\Gamma(U,\mathcal{O}_U)$, $\varphi$ étant l'injection
canonique (\hyperref[I.8.2.7-fr]{I, 8.2.7}). Comme le morphisme $h$ vérifie
les hypothèses de (\hyperref[II.7.3.11-fr]{7.3.11}),
$\Gamma(U,\mathcal{O}_U)$ est un sous-anneau de la fermeture intégrale de $A$
dans $K$, donc est identique à $A$ par hypothèse.
\end{proof}

\begin{remark}[8.11.3]
\label{II.8.11.3-fr}
On verra au chapitre III (\hyperref[III.4.4.11-fr]{III, 4.4.11}) que lorsque
$V$ est un préschéma \emph{localement noethérien}, tout morphisme $h:U\to V$
qui est propre et quasi-fini (en particulier tout morphisme vérifiant les
hypothèses de (\hyperref[II.8.11.2-fr]{8.11.2})) est nécessairement
\emph{fini}. La conclusion de (\hyperref[II.8.11.2-fr]{8.11.2}) résulte donc
dans ce cas de (\hyperref[II.6.1.15-fr]{6.1.15}).
\end{remark}

\begin{env}[8.11.4]
\label{II.8.11.4-fr}
Nous allons maintenant voir que, dans le critère de Grauert
(\hyperref[II.8.9.1-fr]{8.9.1}), on peut souvent affirmer que le préschéma
$C$ et la «~contraction~» $q$ sont déterminés de façon essentiellement
unique.
\end{env}

\begin{lemma}[8.11.5]
\label{II.8.11.5-fr}
Soient $Y$ un préschéma, $p:X\to Y$ un morphisme propre, $\mathcal{L}$ un
$\mathcal{O}_X$-Module inversible $p$-ample, $C$ un $Y$-préschéma,
$\varepsilon:Y\to C$ une $Y$-section,
$q:V=\mathbf{V}(\mathcal{L})\to C$ un $Y$-morphisme, tels que le diagramme
(\hyperref[II.8.9.1.1-fr]{8.9.1.1}) soit commutatif. On suppose en outre que,
si $p=(p_0,\theta)$,
$\theta:\mathcal{O}_Y\to p_*(\mathcal{O}_X)$ est un isomorphisme. Soient alors
$\mathcal{S}'=\bigoplus_{n\geqslant0}p_*(\mathcal{L}^{\otimes n})$,
$C'=\operatorname{Spec}(\mathcal{S}')$, et
$q':\mathbf{V}(\mathcal{L})\to C'$ le $Y$-morphisme canonique
(\hyperref[II.8.8.5-fr]{8.8.5}). Il existe un $Y$-morphisme $u:C'\to C$ et
un seul tel que $q=u\circ q'$.
\end{lemma}

\begin{proof}
L'hypothèse sur $\theta$ entraîne en particulier que $p$ est surjectif~;
comme en vertu de (\hyperref[II.8.8.4-fr]{8.8.4}) la restriction de $q'$ à
$\mathbf{V}(\mathcal{L})-j(X)$ est un \emph{isomorphisme} sur
$C'-\varepsilon'(Y)$ ($\varepsilon'$ étant la section sommet de $C'$), il
résulte de (\hyperref[II.8.8.4-fr]{8.8.4}) que $q'$ est \emph{propre et
surjectif}~; en outre, en vertu de (\hyperref[II.8.8.6-fr]{8.8.6}), si on pose
$q'=(q'_0,\tau)$,
$\tau:\mathcal{O}_{C'}\to q'_*(\mathcal{O}_V)$ est un isomorphisme. On est
donc dans les conditions d'application du lemme
(\hyperref[II.8.11.1-fr]{8.11.1}), et le lemme sera démontré si l'on prouve
que $q$ est constant sur toute fibre $q'{}^{-1}(z')$, où $z'\in C'$. Or, la
condition est trivialement vérifiée pour $z'\notin\varepsilon'(Y)$. D'autre
part, si $z'\in\varepsilon'(Y)$, il existe un $y\in Y$ et un seul tel que
$z'=\varepsilon'(y)$~; et en vertu de la commutativité de
(\hyperref[II.8.8.5.2-fr]{8.8.5.2}) et le fait que $q'$ applique
$\mathbf{V}(\mathcal{L})-j(X)$ dans $C'-\varepsilon'(Y)$,
$q'{}^{-1}(z')=j(p^{-1}(y))$~; la commutativité du diagramme
(\hyperref[II.8.9.1.1-fr]{8.9.1.1}) établit donc notre assertion.
\end{proof}

\begin{corollary}[8.11.6]
\label{II.8.11.6-fr}
Sous les hypothèses de (\hyperref[II.8.11.5-fr]{8.11.5}), supposons outre que
$q$ soit propre et que la restriction de $q$ à
$\mathbf{V}(\mathcal{L})-j(X)$ soit un isomorphisme sur
$C-\varepsilon(Y)$. Alors le morphisme $u$ est universellement fermé,
surjectif et radiciel, et sa restriction à $C'-\varepsilon'(Y)$ est un
isomorphisme sur $C-\varepsilon(Y)$.
\end{corollary}

\begin{proof}
Comme $q'$ est un isomorphisme de $\mathbf{V}(\mathcal{L})-j(X)$ sur
$C'-\varepsilon'(Y)$ (\hyperref[II.8.8.4-fr]{8.8.4}), la dernière assertion
résulte aussitôt de $q=u\circ q'$. En outre, la commutativité des diagrammes

